%% vim:tw=66:spell:wrap:ft=tex:
% Preamble
\documentclass[%
        hyperref={%
                pdfauthor={Zakariyya Mughal},%
                pdfpagemode={None},pdfpagelayout={SinglePage}},%
        xcolor={x11names}%
]{beamer}
\ifx \printpresentnotes \undefined
% no notes
\else
\setbeameroption{show only notes}
\fi

%\usepackage[small,sf,bf]{titlesec}
\PassOptionsToPackage{table,x11names,dvipsnames,rgb}{xcolor}

\usepackage[utf8]{inputenc}
\usepackage[british]{babel}

\usepackage{xstring}

\usepackage{tikz}
\usetikzlibrary{arrows,shapes,automata,positioning,trees,shadows,mindmap,decorations}
%% From <http://tex.stackexchange.com/questions/137357/how-to-draw-an-arrow-with-two-colors>
%% <http://tex.stackexchange.com/a/137438>
\usetikzlibrary{fadings,decorations.pathmorphing}

\makeatletter
\newif\iftikz@shading@path

\tikzset{
    % There are three circumstances in which the fading sep is needed:
    % 1. Arrows which do not update the bounding box (which is most of them).
    % 2. Line caps/joins and mitres that extend outside the natural bounding 
    %    box of the path (these are not calculated by PGF).
    % 3. Other reasons that haven't been anticipated.
    shading xsep/.store in=\tikz@pathshadingxsep,
    shading ysep/.store in=\tikz@pathshadingysep,
    shading sep/.style={shading xsep=#1, shading ysep=#1},
    shading sep=0.0cm,
}

\def\tikz@shadepath#1{% 
    % \tikz@addmode installs the `modes' (e.g., fill, draw, shade) 
    % to be applied to the path. It isn't usualy for doing more
    % changes to the path's construction.
    \iftikz@shading@path%
    \else%
        \tikz@shading@pathtrue%
        % Get the current path.
        \pgfgetpath\tikz@currentshadingpath%
        % Get the shading sep without setting any other keys.
        \begingroup%
            \pgfsys@beginscope% <- may not be necessary
            \tikzset{#1}%
            \xdef\tikz@tmp{\noexpand\def\noexpand\tikz@pathshadingxsep{\tikz@pathshadingxsep}%
                \noexpand\def\noexpand\tikz@pathshadingysep{\tikz@pathshadingysep}}%
            \pgfsys@endscope%
        \endgroup
        \tikz@tmp%
        % Get the boudning box of the current path size including the shading sep
        \pgfextract@process\pgf@shadingpath@southwest{\pgfpointadd{\pgfqpoint{\pgf@pathminx}{\pgf@pathminy}}%
            {\pgfpoint{-\tikz@pathshadingxsep}{-\tikz@pathshadingysep}}}%%
        \pgfextract@process\pgf@shadingpath@northeast{\pgfpointadd{\pgfqpoint{\pgf@pathmaxx}{\pgf@pathmaxy}}%
            {\pgfpoint{\tikz@pathshadingxsep}{\tikz@pathshadingysep}}}%
        % Clear the path
        \pgfsetpath\pgfutil@empty%                          
        % Save the current drawing mode and options.
        \let\tikz@options@saved=\tikz@options%
        \let\tikz@mode@saved=\tikz@mode%
        \let\tikz@options=\pgfutil@empty%
        \let\tikz@mode=\pgfutil@empty%
        % \tikz@options are processed later on.
        \tikz@addoption{%
            \pgfinterruptpath%
            \pgfinterruptpicture%
                \begin{tikzfadingfrompicture}[name=.]
                \pgfscope%
                    \tikzset{shade path/.style=}% Make absolutely sure shade path is not inherited.
                    \path \pgfextra{%
                        % Set the softpath. Any transformations,draw=none} in #1 will have no effect.
                        % This will *not* update the bounding box...
                        \pgfsetpath\tikz@currentshadingpath%
                        % ...so it is done manually.
                        \pgf@shadingpath@southwest
                        \expandafter\pgf@protocolsizes{\the\pgf@x}{\the\pgf@y}%
                        \pgf@shadingpath@northeast%
                        \expandafter\pgf@protocolsizes{\the\pgf@x}{\the\pgf@y}%
                        % Install the drawing modes and options.
                        \let\tikz@options=\tikz@options@saved%
                        \let\tikz@mode=\tikz@mode@saved%
                    };
                    % Now get the bounding box of the picture.
                    \xdef\pgf@shadingboundingbox@southwest{\noexpand\pgfqpoint{\the\pgf@picminx}{\the\pgf@picminy}}%
                    \xdef\pgf@shadingboundingbox@northeast{\noexpand\pgfqpoint{\the\pgf@picmaxx}{\the\pgf@picmaxy}}%
                    \endpgfscope
                \end{tikzfadingfrompicture}%
            \endpgfinterruptpicture%
            \endpgfinterruptpath%
            % Install a rectangle that covers the shaded/faded path picture.
            \pgftransformreset%
            \pgfpathrectanglecorners{\pgf@shadingboundingbox@southwest}{\pgf@shadingboundingbox@northeast}%
            %
            % Reset all modes.
            \let\tikz@path@picture=\pgfutil@empty%
            \tikz@mode@fillfalse%
            \tikz@mode@drawfalse%
            \tikz@mode@tipsfalse%
            \tikz@mode@doublefalse%
            \tikz@mode@clipfalse%
            \tikz@mode@boundaryfalse%
            \tikz@mode@fade@pathfalse%
            \tikz@mode@fade@scopefalse%
            % Now install shading options.
            \tikzset{#1}%
            \tikz@mode%
            % Make the fading happen.
            \def\tikz@path@fading{.}%
            \tikz@mode@fade@pathtrue%
            \tikz@fade@adjustfalse%
            % Shift the fading to the mid point of the rectangle
            \pgfpointscale{0.5}{\pgfpointadd{\pgf@shadingboundingbox@southwest}{\pgf@shadingboundingbox@northeast}}%
            \edef\tikz@fade@transform{shift={(\the\pgf@x,\the\pgf@y)}}%
            \pgfsetfading{\tikz@path@fading}{\tikz@do@fade@transform}%
            \tikz@mode@fade@pathfalse%              
        }%
    \fi%
}
\tikzset{
    shade path/.code={%
        \tikz@addmode{\tikz@shadepath{#1}}%
    }
}

% requires tikz
\usepackage{colortbl}

\usepackage{float}
\usepackage{graphicx}
\usepackage[normalem]{ulem}

%% use as
%%     \vertcenterimage{\includegraphics{*}}
\newcommand{\vertcenterimage}[1]{\raisebox{-.5\height}{#1}}

%% use as
%%     \flipbox{\includegraphics{*}}
\newcommand{\flipbox}[1]{\scalebox{1}[-1]{#1}}

\usepackage{amsmath}
\usepackage{amsfonts}
\usepackage{bm} % bold mathematics (\bm command)

% set-builder notation using \Set{ ... | ... }
\usepackage{braket}

\usepackage{mathtools}
\usepackage{breqn}

%% make everything in the box bold; both the text and mathematics: using
%% \boldmath from `bm` package
\newcommand{\be}{\bfseries\boldmath}

\usepackage{pifont}% http://ctan.org/pkg/pifont
\newcommand{\cmark}{\ding{51}}%
\newcommand{\xmark}{\ding{55}}%

\newcommand{\todo}[1]{%
	\textcolor{red}{TODO: #1}%
}
\newcommand{\todofig}[1]{%
	\textcolor{red}{TODO figure: \nolinkurl{#1}}%
}


\usepackage{hyperref}
\hypersetup{%
	pdfauthor={Zakariyya Mughal},%
	pdfpagemode={None},%
	pdfpagelayout={SinglePage}%
}


\usetheme{Ilmenau}
\usecolortheme{beaver}
\usepackage{ifxetex}

\ifxetex
	% set font to Tahoma
	\usefonttheme{professionalfonts} % using non standard fonts for beamer
	\usefonttheme{serif} % default family is serif
	\usepackage{fontspec}
	\setmainfont{Tahoma}
\else
	\usepackage[T1]{fontenc}
\fi

\definecolor{CBLRed}{RGB}{205,0,0}

\colorlet{langC}{blue!40}
\colorlet{langM}{orange!20}
\colorlet{data}{grey!10}

% Title on title slide
\setbeamerfont{title}{size = \Large}
\setbeamercolor{title}{fg = black, bg = white}

\setbeamercolor{frametitle}{fg=white,bg=CBLRed}
\setbeamercolor{section in head/foot}{fg=white,bg=CBLRed}
\setbeamercolor{subsection in head/foot}{fg=white,bg=CBLRed}
\setbeamercolor{author in head/foot}{fg=white,bg=CBLRed}
\setbeamertemplate{headline}{}
\beamertemplatenavigationsymbolsempty

\logo{}
\def\insertlogo
{%
	\color{gray}
	\begin{tabular*}{\paperwidth}{p{0.01\paperwidth}p{0.84\paperwidth}p{0.15\paperwidth}}
		&
		\parbox[c][][c]{\textwidth}{\raggedright
			\vertcenterimage{\includegraphics[width=\textwidth]{gfx/cbl-footer.png}}
		}
		&
		\parbox[c][][c]{\textwidth}{%\raggedleft
			{\large\insertframenumber} %/\inserttotalframenumber
		}
	\end{tabular*}
 }

%\usefonttheme[onlymath]{serif}
\ifx \printpresentarticle \undefined
	\setbeamertemplate{frametitle}[default][center]
	\setbeamertemplate{footline}{}
	%\setbeamertemplate{footline}[frame number]

	%% Change beamer bullets to circles rather than the ball default
	\setbeamertemplate{itemize items}[circle]
	\setbeamertemplate{enumerate items}[circle]
\fi


\usepackage{textcomp}
\usepackage{fancyvrb}
\usepackage{changepage}
\usepackage{multicol}
\usepackage{wasysym}
\usepackage{listings}

\lstset{%basicstyle=\small\ttfamily,
%numbers=left,
%escapeinside=||
}
\newenvironment{indented}{\begin{adjustwidth}{1.5em}{}}{\end{adjustwidth}}

% http://tex.stackexchange.com/questions/12550/changing-default-width-of-blocks-in-beamer/12551#12551
\newenvironment<>{varblock}[2][.9\textwidth]{%
  \setlength{\textwidth}{#1}
  \begin{actionenv}#3%
    \def\insertblocktitle{#2}%
    \par%
    \usebeamertemplate{block begin}}
  {\par%
    \usebeamertemplate{block end}%
  \end{actionenv}}

\ifx \printpresentnote \undefined
% no notes
\else
\setbeameroption{show only notes}
\fi

%% use as
%%     \af{ number of overlay }{ slide label }
%% e.g.,
%%     \af{2}{intro-slide}
\newcommand{\af}[2]{\againframe<#1|handout:0>[noframenumbering]{#2}}


%% TOC at beginning of sections/subsections
\ifx \printpresenthandout \undefined
	% no handout
	\AtBeginSubsection[]
	{
		\begin{frame}[noframenumbering]\frametitle{Table of contents}
			\begin{multicols}{2}
				\tableofcontents[
					currentsection,
					hideothersubsections,
					sectionstyle=show/shaded,
					subsectionstyle=show/shaded/hide,
				]
			\end{multicols}
	  \end{frame}
	}
\else
	\AtBeginSection[]
	{
		\begin{frame}[noframenumbering]\frametitle{Table of contents}
			\begin{multicols}{2}
				\tableofcontents[
					currentsection,
					currentsubsection,
					hideothersubsections,
					sectionstyle=show/shaded,
					subsectionstyle=show/show/hide,
				]
			\end{multicols}
	  \end{frame}
	}
\fi

%% TikZ arrows
%% From <https://tex.stackexchange.com/questions/61507/drawing-arrows-in-beamer>
\tikzset{
    myarrow/.style={
        draw,
        fill=orange,
        single arrow,
        minimum height=3.5ex,
        single arrow head extend=1ex
    }
}
\newcommand{\arrowup}{%
\tikz [baseline=-0.5ex]{\node [myarrow,rotate=90] {};}
}
\newcommand{\arrowdown}{%
\tikz [baseline=-1ex]{\node [myarrow,rotate=-90] {};}
}
\newcommand{\arrowright}{%
\tikz [baseline=-0.5ex]{\node [myarrow,rotate=0] {};}
}
\newcommand{\arrowleft}{%
\tikz [baseline=-0.5ex]{\node [myarrow,rotate=180] {};}
}


\definecolor{quotemark}{gray}{0.7}
\newenvironment{fancyquote}%
	{%
	    \vspace{1em}%
	    \singlespacing
	    \noindent%
		 \begin{picture}(0,0)%
		 \put(-15,-0){\makebox(0,0){\scalebox{3}{\textcolor{quotemark}{``}}}}%
		 \end{picture}%
	\footnotesize\upshape%
	}%
	{%
	 \par%
	 \makebox[0pt][l]{%
	 \hspace{\linewidth}%
	 \begin{picture}(0,0)(0,0)%
	 \put(15,20){\makebox(0,0){%
	 \scalebox{3}{\color{quotemark}''}}}%
	 \end{picture}}%
	   \vspace{-2.5em}%
	}%


\usepackage{diagbox}

\begin{document}
% meta needs to be in \begin{document} so that tabular works
\title[]{Converting a Neuron Morphology Reconstruction System: Open Science Design and Implementation}
\author[]{Zakariyya Mughal}
\institute[]{Computational Biomedicine Lab\\University of Houston}
\date{2015 Nov 30}

%  - Title slide.%{{{
\frame{\titlepage}
%}}}


\section{Introduction}
%  - Motivation - short introduction to your research, pointers to applications, why the research is important.
\subsection{Motivation}
\begin{frame}
\frametitle{Motivation}
\end{frame}

%  - Problem Statement - Goals, objectives, aims. May also include challenges.
%        IMPORTANT: State your specific aims clearly - what you plan to achieve and how
%        you plan to achieve it. Let your contributions to the field of Computer Science
%        be clear.
\subsection{Problem}
\begin{frame}<1>[label=prob]\frametitle{Problems}
\begin{itemize}
	\item<1-> Neuron reconstruction
	\item<3-> Open science
	\item<2-> BigNeuron
\end{itemize}
\end{frame}

\begin{frame}\frametitle{Neuron reconstruction}
	% TODO
	[image of neuron volume]
	$\rightarrow$
	[image of neuron tracing]
	\note{
		problem with manual tracing
	}
\end{frame}

\againframe<2>{prob}

\begin{frame}\frametitle{Open science}
	% TODO
	Human Genome Project
	% First example of open research
	% <https://en.wikipedia.org/wiki/Bermuda_Principles> ---> 1996
	% - data available in 24 hours of sequence assembly
	% - immediate publication of annotated sequences
	% - all data in the public domain
	Allen Brain Atlas
	OpenWorm
	Galaxy Zoo
\end{frame}

\againframe<3>{prob}

\begin{frame}\frametitle{BigNeuron}
	% TODO
	[Vaa3D GUI image]
\end{frame}

\begin{frame}[label=challenges]
\frametitle{Challenges and risks}
\begin{itemize}
	\item Starting without a test suite
	\item Memory management
	\item Data layout differences of n-dimensional arrays
	\item Data handling differences
	\item Reimplement MATLAB Toolbox functions
\end{itemize}
\end{frame}

\subsection{Objectives}
\begin{frame}
\frametitle{Objectives}
\begin{itemize}
	\item Integrate with Vaa3D tool for biomedical image analysis
	\item Convert the MATLAB code to C / C++
	\item Release code as free/open source software for open-science
	\item Build a test suite
\end{itemize}
\end{frame}

\section{Analysis}

\section{Design}

\begin{frame}
\frametitle{Design}
\begin{itemize}
	\item Use the original MATLAB code as initial design: followed by refactoring the C code
	\item Automated GNU Make-based build system:
		\begin{itemize}
			\item compiler extension configuration
			\item third-party library build flag configuration
			\item prerequisite scanning
		\end{itemize}
	\item Available on GitHub for public visibility
\end{itemize}
\end{frame}

\section{Implementation}

\begin{frame}
\frametitle{Implementation}
\begin{itemize}
	\item Base data structure is an n-dimensional array (row-major buffer with shape metadata)
	\item float32 for storage; float64 for calculations
	\item In-memory calculations
	\item ITK library used for implementing several filters
	\item Branch-based component development
\end{itemize}
\end{frame}

\section{Testing}

\begin{frame}
\frametitle{Testing}
\begin{itemize}
	\item Use the Test Anything Protocol (TAP) for test output
	\item Continuous integration with a pull request based workflow
	\item Automatic provisioning of build environment
	\item Tracing-based comparison to original MATLAB function
	\item Assertions for pre-/post-conditions
	\item External dependencies managed using version control
	\item Integration testing with Vaa3D plugin
\end{itemize}
\end{frame}

\section{Discussion}

\section{Conclusions}

\begin{frame}
\frametitle{Conclusions}
\begin{itemize}
	\item Software is an ever-changing system
	\item The code depends on an ecosystem
	\item Rewrites across languages are difficult
	\item Scientific software engineering is very much like traditional software engineering
	\item Open-science and reproducibility for computational methods requires software engineering skills
\end{itemize}
\end{frame}

%  - Related Work - a catchy name is “State of the Art”.
%        Literature review relevant to specific research. Your committee
%        members are expected to already have a general background in your field.
%        Concentrate on the details of your specific research area. It is preferable to
%        use more recent literature (except for some classical “old” ones). Also avoid
%        reviewing standard concepts that have been around for a long time.
%  - Proposed Framework - this is the most important part of your proposal.
%        Expect most questions from committee members on this. Put in all your
%        best here. Describe the entire framework before separating into work-done and
%        work-left. Do not mix framework with results. You may start with a one-slide
%        overview of the entire framework before going into details. Make sure committee
%        members will be able to understand your work. Relate to previous approaches and
%        how your is an extension or a difference. Include mathematical equations that
%        will explain your work.
%  - Explain the relevance of each equation. Understand the meaning and
%        significance of each term in the equation irrespective of whether you define
%        them in the slide or not. TIP: Rather than doing your equation twice (for your
%        PPT and Latex write-up), you may do it once in your Latex document and copy
%        into TexPoint for the PPT.
%  - Preliminary Results - if you already have some results, please show it.
%        Talk about how you propose to validate your results. Show how your
%        results compare with previous results OR how it performs in general, (e.g., in
%        how many images it succeeded or in how many it failed). A chart will be useful.
%        Any statistical comparison should indicate level of significance.
%  - Present status of the proposed framework
%        - which aspects are completed, which are in progress, and which have not yet been done at all.
%  - It is advisable to include a timeline of future indicating milestones and when you hope to achieve each.
%  - Expected Impact - mention specific application areas where you have tried or will try your proposed framework.
%        Also, explain how it can be generalized to other areas in Computer
%        Science. What is the overall benefit for the community? What is the
%        contribution to Computer Science research?
%  - Conclusion slide could be catchy, something they will remember.
%        Maybe an animation, video, or terse statement that summarizes it all.
%  - Publications - you should list your publications or mention them in relevant sections of the presentation. 

\end{document}
